\chapter{A Production Path Tracer: Media, Randomness and Validation}\label{03-visual-computing-i-ch010-a-production-path-tracer-media-randomness-and-validation}\phantomsection\hypertarget{03-visual-computing-i-ch010}{}\label{stagechapter:03-visual-computing-i-ch010}

\hypertarget{03-visual-computing-i-ch010-a-production-path-tracer-media-randomness-and-validation}{}

\hyperref[volume-reading-map]{Reading map} · \hyperref[03-visual-computing-i-ch009-point-lights-mis-and-environment-lighting]{Previous chapter} · \href{../06-GPU-Architecture-and-Real-Time-Rendering/main.pdf\#nameddest=03-visual-computing-i-ch011-bvh-traversal-and-deterministic-evidence}{Next chapter (GPU Architecture and Real-Time Rendering, Ch. 8)}

A production path tracer must handle nested media, random streams, execution and regression evidence as well as the core estimator. This span joins those engineering contracts to implementation and profiling.

\textbf{By the end:} Identify the implementation and validation contracts needed for a robust path tracer.

\textbf{Recall question:} Which failures require a mathematical oracle, and which require scene-level tests?

\subsection{Next --- Pass 8: M125 + M126, Parts 2070--2100}\label{03-visual-computing-i-ch010-next-pass-8-m125-m126-parts-20702100}

\hypertarget{03-visual-computing-i-ch010-next-pass-8-m125-m126-parts-20702100}{}

That is the \textbf{production path-tracer completion pass}: nested dielectric media and IOR stacks, TIR in the path integrator, colored-throughput Russian roulette, deterministic/stratified/low-discrepancy/blue-noise streams, multithreaded tile scheduling, per-thread RNGs, false-sharing-safe film accumulation, HDR output, debug passes, convergence/reference tests, rays/sec and BVH profiling, packet/wavefront/GPU ray architectures, then a complete end-to-end path tracer with fault isolation across geometry, PDF, BRDF, visibility and accumulation.

Pass 8 completes \textbf{Phase 21 --- Ray Tracing, Monte Carlo Light Transport, Sampling, BVH Rendering, and Path Tracing}.

The curriculum explicitly positions M125 as the ``make transport survive real materials, real randomness and real CPU execution'' module and M126 as the ``prove it, profile it, restructure it for parallel hardware, and know exactly where it is broken'' module.

Pass 7 gave us:

\[
\beta_{k+1}
=
\beta_k
\frac{
f_k|\cos\theta_k|
}{
p_k
}
\]

and:

\[
L
\leftarrow
L+
\beta L_e
\]

plus NEE and MIS.

Now we make that system survive scenes like:

\begin{verbatim}
camera
  ↓
air
  ↓
water
  ↓
glass shell
  ↓
air bubble
  ↓
glass
  ↓
water
\end{verbatim}

while hundreds of millions of rays are generated by many CPU/GPU workers without races, repeated random patterns, broken accumulation, incorrect PDFs, or untraceable rendering errors.

\begin{center}\rule{0.5\linewidth}{0.5pt}\end{center}

\subsection{PARTS 2070--2089}\label{03-visual-computing-i-ch010-parts-20702089}

\hypertarget{03-visual-computing-i-ch010-parts-20702089}{}

\textbf{Why simple glass breaks the moment media are nested}

The Whitted renderer we inspected earlier essentially assumes:

\[
\text{outside medium}=\text{air},
\]

and:

\[
\text{inside medium}=\text{material.IOR}.
\]

So when entering glass:

\[
\eta_i=1,
\qquad
\eta_t=1.5.
\]

When leaving:

\[
\eta_i=1.5,
\qquad
\eta_t=1.
\]

That works for:

\begin{verbatim}
air → isolated glass object → air.
\end{verbatim}

But now place the glass object under water.

Correct transition is:

\[
\text{water}
\rightarrow
\text{glass}.
\]

Therefore:

\[
\eta_i\approx1.333,
\]

\[
\eta_t\approx1.5.
\]

If the renderer still uses:

\[
1\rightarrow1.5,
\]

Snell's law is wrong.

So are:

\begin{itemize}
\tightlist
\item
  refraction direction,
\item
  Fresnel reflectance,
\item
  total-internal-reflection conditions.
\end{itemize}

The visible result can be substantially wrong even though the glass shader ``looks plausible.''

\begin{center}\rule{0.5\linewidth}{0.5pt}\end{center}

\textbf{Refraction depends on a pair of media}

Snell's law:

\[
\boxed{
\eta_i\sin\theta_i
=
\eta_t\sin\theta_t.
}
\]

The quantity needed in the vector refraction equation is:

\[
\boxed{
\eta=
\frac{\eta_i}{\eta_t}.
}
\]

The important conceptual shift is:

\[
\boxed{
\text{IOR is not solely a property of the surface interaction.}
}
\]

It describes a \textbf{medium}.

The interface tells us which medium exists on each side.

A physically useful surface boundary therefore conceptually contains something like:

\begin{verbatim}
struct MediumInterface
{
    MediumID Outside;
    MediumID Inside;
};
\end{verbatim}

rather than merely:

\begin{verbatim}
float MaterialIOR;
\end{verbatim}

\begin{center}\rule{0.5\linewidth}{0.5pt}\end{center}

\textbf{Properly nested dielectric media}

Consider:

\[
\text{air}=1.0003,
\]

\[
\text{water}=1.333,
\]

\[
\text{glass}=1.5.
\]

A camera ray starts in air.

Medium state:

\begin{verbatim}
[air]
\end{verbatim}

It enters water:

\begin{verbatim}
[air, water]
\end{verbatim}

Transition:

\[
1.0\rightarrow1.333.
\]

It enters a glass object submerged in the water:

\begin{verbatim}
[air, water, glass]
\end{verbatim}

Transition:

\[
1.333\rightarrow1.5.
\]

It exits glass back into water:

\begin{verbatim}
[air, water]
\end{verbatim}

Transition:

\[
1.5\rightarrow1.333.
\]

Then leaves water:

\begin{verbatim}
[air]
\end{verbatim}

Transition:

\[
1.333\rightarrow1.0.
\]

That is the intuition behind an \textbf{IOR stack}.

\begin{center}\rule{0.5\linewidth}{0.5pt}\end{center}

\textbf{Why a stack works only for properly nested closed volumes}

A simple stack assumes interfaces obey properly nested topology.

For example:

\begin{verbatim}
air
 └─ water
     └─ glass
         └─ air bubble
\end{verbatim}

Every entry has a corresponding exit in reverse order.

Then push/pop is well defined.

But arbitrary overlapping meshes can look like:

\begin{verbatim}
object A enters
object B enters
object A exits
object B exits
\end{verbatim}

which is not a proper stack nesting.

Production renderers therefore often distinguish:

\[
\text{geometric object identity}
\]

from:

\[
\text{medium interface identity}
\]

and may require explicit medium priorities or robust boundary semantics.

For our Kairo progression, a stack is still the correct next architecture because it handles the common closed-nested case cleanly.

\begin{center}\rule{0.5\linewidth}{0.5pt}\end{center}

\textbf{A path needs medium state}

A production path state can contain:

\begin{verbatim}
struct MediumStack
{
    std::array<MediumID, 8> IDs{};
    uint32_t Count = 1;

    MediumID Current() const
    {
        return IDs[Count - 1];
    }
};
\end{verbatim}

Start with:

\begin{verbatim}
IDs[0] = Air;
\end{verbatim}

When a transmission event enters a dielectric volume:

\begin{verbatim}
push(interiorMedium);
\end{verbatim}

When it exits:

\begin{verbatim}
pop();
\end{verbatim}

Reflection does \textbf{not} change the medium.

This distinction is essential.

\begin{center}\rule{0.5\linewidth}{0.5pt}\end{center}

\textbf{Front-face information becomes mandatory}

We previously noted that Kairo's current \texttt{SurfaceHit} stores:

\begin{verbatim}
Position
Normal
UV
PrimitiveIndex
MaterialIndex
\end{verbatim}

but does not preserve separate:

\[
N_g,
\]

\[
N_s,
\]

or:

\[
\text{FrontFace}.
\]

That becomes a real architectural limitation now.

To update a medium stack correctly, the path tracer must know whether the ray crossed:

\[
\text{outside}\rightarrow\text{inside}
\]

or:

\[
\text{inside}\rightarrow\text{outside}.
\]

Use the geometric normal:

\[
\boxed{
\text{frontFace}
=
D\cdot N_g<0.
}
\]

Not the normal-mapped shading normal.

A normal map must never cause the renderer to suddenly believe the ray entered a different dielectric volume.

\begin{center}\rule{0.5\linewidth}{0.5pt}\end{center}

\textbf{Upgraded hit record}

A production-oriented Kairo hit would benefit from:

\begin{verbatim}
struct SurfaceHit
{
    float T = Infinity;

    Vec3 Position;

    // True primitive orientation.
    Vec3 GeometricNormal;

    // Smooth/normal-mapped normal used for BSDF evaluation.
    Vec3 ShadingNormal;

    Vec2 UV;
    Vec3 Barycentric;

    uint32_t PrimitiveIndex;
    uint32_t MaterialIndex;

    bool FrontFace;
};
\end{verbatim}

Maintain:

\[
N_s\cdot N_g\ge0.
\]

Use:

\[
N_g
\]

for:

\begin{itemize}
\tightlist
\item
  entering/exiting classification,
\item
  origin offset,
\item
  medium crossing.
\end{itemize}

Use:

\[
N_s
\]

for most BRDF calculations.

That separation prevents an entire class of path-tracing bugs.

\begin{center}\rule{0.5\linewidth}{0.5pt}\end{center}

\textbf{Computing \(\eta_i\) and \(\eta_t\)}

Suppose current medium stack is:

\begin{verbatim}
[air, water]
\end{verbatim}

and hit is front-facing glass boundary.

Then:

\[
\eta_i=
\eta(\text{water})
=
1.333.
\]

Target interior medium:

\[
\eta_t=
1.5.
\]

If the ray instead reaches the back face of that glass shell, current medium is glass:

\[
\eta_i=1.5
\]

and after leaving it the stack underneath says:

\[
\eta_t=1.333.
\]

The previous medium, not automatically air, is the target.

\begin{center}\rule{0.5\linewidth}{0.5pt}\end{center}

\textbf{Refraction with tracked media}

Let incoming travel direction be:

\[
I.
\]

Choose geometric interface normal \(N\) to oppose the incoming ray.

Then:

\[
\cos\theta_i
=
-N\cdot I.
\]

Ratio:

\[
\eta=
\frac{\eta_i}{\eta_t}.
\]

Compute:

\[
\sin^2\theta_t
=
\eta^2
\left(
1-\cos^2\theta_i
\right).
\]

If:

\[
\sin^2\theta_t\le1,
\]

then:

\[
\cos\theta_t
=
\sqrt{
1-\sin^2\theta_t
}.
\]

Transmission direction:

\[
\boxed{
T
=
\eta I+
(
\eta\cos\theta_i-\cos\theta_t
)N.
}
\]

\begin{center}\rule{0.5\linewidth}{0.5pt}\end{center}

\textbf{Total internal reflection}

If:

\[
\eta_i>\eta_t,
\]

a sufficiently grazing incident ray cannot refract.

Condition:

\[
\boxed{
\sin^2\theta_t>1.
}
\]

Then no real \(\cos\theta_t\) exists.

The correct path event becomes:

\[
\boxed{
\text{reflection probability}=1.
}
\]

No transmission event occurs.

And crucially:

\[
\boxed{
\text{the medium stack must not change.}
}
\]

The ray did not cross the interface.

\begin{center}\rule{0.5\linewidth}{0.5pt}\end{center}

\textbf{Critical angle}

From Snell:

\[
\eta_i\sin\theta_c
=
\eta_t.
\]

Therefore:

\[
\boxed{
\theta_c
=
\arcsin
\left(
\frac{\eta_t}{\eta_i}
\right)
}
\]

when:

\[
\eta_i>\eta_t.
\]

For glass:

\[
\eta_i=1.5,
\]

to air:

\[
\eta_t=1,
\]

we get:

\[
\theta_c
=
\arcsin(2/3)
\approx41.81^\circ.
\]

For glass to water:

\[
\theta_c
=
\arcsin
\left(
\frac{1.333}{1.5}
\right)
\approx62.7^\circ.
\]

Notice the enormous difference.

A renderer that assumes every glass exit goes to air will produce TIR where real submerged glass should transmit.

That is not a subtle error.

\begin{center}\rule{0.5\linewidth}{0.5pt}\end{center}

\textbf{Fresnel must use the correct medium pair too}

At normal incidence, dielectric Fresnel:

\[
\boxed{
F_0
=
\left(
\frac{
\eta_i-\eta_t
}{
\eta_i+\eta_t
}
\right)^2.
}
\]

Air to glass:

\[
F_0
=
\left(
\frac{1-1.5}{1+1.5}
\right)^2
=
0.04.
\]

Water to glass:

\[
F_0
=
\left(
\frac{1.333-1.5}{1.333+1.5}
\right)^2
\approx0.00347.
\]

Only about:

\[
0.35\%.
\]

That is radically different from 4\%.

This explains why submerged interfaces can almost disappear visually.

Using fixed:

\[
F_0=0.04
\]

for every dielectric interface is only valid for approximately air-to-\(1.5\)-IOR surfaces.

\begin{center}\rule{0.5\linewidth}{0.5pt}\end{center}

\textbf{Glass event selection with medium state}

A path tracer should calculate:

\[
F(\eta_i,\eta_t,\theta_i).
\]

Then:

\[
P(\text{reflection})=F,
\]

\[
P(\text{transmission})=1-F.
\]

If TIR:

\[
F=1.
\]

If reflection chosen:

medium stack unchanged.

If transmission chosen:

medium state crosses the boundary.

That order matters.

You should not push/pop the medium before knowing whether the path actually transmitted.

\begin{center}\rule{0.5\linewidth}{0.5pt}\end{center}

\textbf{A useful medium-aware dielectric sample}

\begin{verbatim}
struct DielectricSample
{
    Vec3 Wi;

    bool Reflection = false;
    bool Transmission = false;
    bool Delta = true;

    float EventProbability = 1.0f;

    // The medium the ray will occupy after transmission.
    MediumID NextMedium;
};

DielectricSample SampleDielectric(
    const SurfaceHit& hit,
    const Ray& ray,
    MediumStack& media,
    const Material& material,
    RNG& rng)
{
    const Vec3 Ng =
        hit.GeometricNormal;

    const bool entering =
        hit.FrontFace;

    // Orient the interface normal against incident travel.
    const Vec3 N =
        entering ? Ng : -Ng;

    const MediumID currentMedium =
        media.Current();

    const MediumID targetMedium =
        entering
            ? material.InteriorMedium
            : media.BelowTop();

    const float etaI =
        GetMedium(currentMedium).IOR;

    const float etaT =
        GetMedium(targetMedium).IOR;

    const float eta =
        etaI / etaT;

    const float cosThetaI =
        std::clamp(
            Dot(-ray.Direction, N),
            0.0f,
            1.0f);

    const float sin2ThetaT =
        eta * eta *
        (1.0f -
         cosThetaI * cosThetaI);

    // Total internal reflection.
    if (sin2ThetaT > 1.0f)
    {
        return
        {
            .Wi =
                Reflect(
                    ray.Direction,
                    N),

            .Reflection = true,
            .EventProbability = 1.0f,
            .NextMedium = currentMedium
        };
    }

    const float F =
        FresnelDielectric(
            cosThetaI,
            etaI,
            etaT);

    if (rng.NextFloat() < F)
    {
        return
        {
            .Wi =
                Reflect(
                    ray.Direction,
                    N),

            .Reflection = true,
            .EventProbability = F,
            .NextMedium = currentMedium
        };
    }

    const Vec3 transmitted =
        Refract(
            ray.Direction,
            N,
            eta);

    return
    {
        .Wi = transmitted,
        .Transmission = true,
        .EventProbability = 1.0f - F,
        .NextMedium = targetMedium
    };
}
\end{verbatim}

Then update the stack only after \texttt{Transmission\ ==\ true}.

\begin{center}\rule{0.5\linewidth}{0.5pt}\end{center}

\textbf{Refractive throughput has an extra subtlety}

For reflection in the same medium, radiance transport is comparatively simple.

Across a refractive interface, solid angle changes according to Snell's law.

Radiance across media obeys an optical invariant involving:

\[
\eta^2.
\]

Consequently a rigorous BSDF implementation must include an appropriate squared-IOR factor for transmission depending on whether it is transporting \textbf{radiance} or its adjoint/importance.

A common camera-to-light radiance-transport convention introduces a factor proportional to:

\[
\boxed{
\left(
\frac{\eta_i}{\eta_t}
\right)^2
}
\]

inside the specular-transmission throughput.

The exact placement depends on your BSDF direction and transport-mode conventions.

The engineering rule is more important than memorizing one isolated expression:

\[
\boxed{
\text{pick one transport convention and apply its } \eta^2 \text{ factor consistently everywhere}.
}
\]

Do not copy a refraction vector from one renderer and a BTDF weight from another renderer whose \(\omega_i/\omega_o\) conventions differ.

\begin{center}\rule{0.5\linewidth}{0.5pt}\end{center}

\textbf{Russian roulette with colored throughput}

Pass 7's generic throughput:

\[
\beta=
(\beta_r,\beta_g,\beta_b).
\]

Suppose:

\[
\beta=
(0.9,0.03,0.01).
\]

Using luminance alone may say this path is relatively dim.

But the red channel remains very important.

A common roulette survival estimate is:

\[
\boxed{
p_s=
\operatorname{clamp}
(
\max(\beta_r,\beta_g,\beta_b),
p_{\min},
p_{\max}
).
}
\]

For the example:

\[
p_s\approx0.9.
\]

On survival:

\[
\boxed{
\beta\leftarrow\frac{\beta}{p_s}.
}
\]

That keeps every color channel unbiased.

\begin{center}\rule{0.5\linewidth}{0.5pt}\end{center}

\textbf{Why roulette cannot divide each channel by a different probability}

For the ordinary shared RGB path described here, use one survival decision and divide every channel by the same survival probability. Independent channel roulette is a different estimator and requires its own compensation and bookkeeping; it is not mathematically impossible.

A path is one geometric transport event.

If red survived while green died, you would conceptually have separate spectral paths with different directions/visibility histories.

RGB path tracing normally keeps one common path event and one common survival probability.

The color remains in:

\[
\beta.
\]

\begin{center}\rule{0.5\linewidth}{0.5pt}\end{center}

\textbf{Refraction and roulette}

An advanced subtlety appears with refractive \(\eta^2\) scaling.

Throughput magnitude can jump because of a change in transport measure rather than because the path has become physically more or less important.

A production renderer may maintain an auxiliary refractive scale when deciding roulette probability so an \(\eta^2\) convention does not cause pathological termination behavior.

The key principle is:

\[
\boxed{
\text{roulette probability should estimate future contribution, not blindly react to representation artifacts}.
}
\]

\begin{center}\rule{0.5\linewidth}{0.5pt}\end{center}

\textbf{Deterministic random-number streams per pixel}

Suppose renderer uses one global RNG protected by a mutex.

Thread A processes pixel 10 first this run.

Thread B processes it first next run.

Random-number consumption order changes.

The image changes.

That makes debugging miserable.

Instead derive random streams from stable identity:

\[
\boxed{
(\text{pixel},\text{sample},\text{dimension},\text{global seed}).
}
\]

Then the same pixel/sample receives the same random sequence regardless of:

\begin{itemize}
\tightlist
\item
  thread count,
\item
  tile order,
\item
  scheduling.
\end{itemize}

Current Kairo already moves strongly in this direction.

Path RNG seed is built from:

\[
x,
y,
sample,
SampleSeed.
\]

So a path's random sequence is tied to the pixel/sample rather than to the worker thread.

That is excellent for reproducibility.

\begin{center}\rule{0.5\linewidth}{0.5pt}\end{center}

\textbf{Current Kairo RNG}

Current seed:

\begin{verbatim}
rng =
    x      * 73856093u ^
    y      * 19349663u ^
    sample * 83492791u ^
    SampleSeed * 2654435761u ^
    0x9e3779b9u;
\end{verbatim}

Then \texttt{TracePath()} advances a local xorshift state.

This gives:

\[
\boxed{
\text{thread-schedule-independent local randomness}.
}
\]

That property is arguably better for deterministic rendering than a pure per-thread stateful RNG.

\begin{center}\rule{0.5\linewidth}{0.5pt}\end{center}

\textbf{The weakness of sequential per-path RNG consumption}

Suppose bounce code currently consumes random dimensions as:

\begin{verbatim}
dimension 0 → BSDF x
dimension 1 → BSDF y
dimension 2 → Russian roulette
\end{verbatim}

Later you add light sampling before BSDF:

\begin{verbatim}
dimension 0 → light x
dimension 1 → light y
dimension 2 → BSDF x
dimension 3 → BSDF y
dimension 4 → roulette
\end{verbatim}

Every later random number changes.

This makes before/after debugging difficult.

A more sophisticated sampler gives named/stable dimensions:

\begin{verbatim}
Sample1D(pixel, sample, Dimension::LightSelection);
Sample2D(pixel, sample, Dimension::LightUV);
Sample2D(pixel, sample, Dimension::BSDF);
Sample1D(pixel, sample, Dimension::Roulette);
\end{verbatim}

Counter-based RNGs and Sobol samplers are especially useful here because a sample can be addressed by:

\[
(\text{sample index},\text{dimension})
\]

without depending on consumption order.

\begin{center}\rule{0.5\linewidth}{0.5pt}\end{center}

\textbf{Correlated random samples}

Two random numbers can each look uniform individually and still be badly correlated together.

Suppose:

\[
\xi_2=\xi_1.
\]

Every 2D sample lies on:

\[
y=x.
\]

Marginals are uniform.

2D distribution is terrible.

For a disk sampler this can produce directional streaks.

For area lights it can produce banding.

For lens samples it can create structured bokeh artifacts.

Rendering needs good \textbf{multidimensional} sample structure, not merely a histogram that looks flat in each dimension.

\begin{center}\rule{0.5\linewidth}{0.5pt}\end{center}

\textbf{Correlation between pixels}

Suppose every pixel uses exactly the same subpixel jitter pattern.

Edges can develop coherent structured artifacts.

Independent randomness converts the structure into noise, but pure white noise is itself not always ideal perceptually.

This is why rendering sampling eventually moves toward:

\[
\text{stratification},
\]

\[
\text{low discrepancy},
\]

\[
\text{blue noise}.
\]

These control where error energy appears.

\begin{center}\rule{0.5\linewidth}{0.5pt}\end{center}

\textbf{Stratified pixel sampling}

Divide pixel into:

\[
n_x\times n_y
\]

cells.

For a \(2\times2\) example:

\begin{verbatim}
+-----+-----+
|  •  | •   |
+-----+-----+
|   • |  •  |
+-----+-----+
\end{verbatim}

Draw one jittered sample per cell.

Compared with four independent random samples, stratification guarantees broad pixel coverage.

For one dimension split into \(N\) strata:

\[
x_i=
\frac{i+\xi_i}{N}.
\]

This reduces clustering.

Current Kairo's \texttt{StratifiedPixelSample()} implements this principle with deterministic hashed jitter.

\begin{center}\rule{0.5\linewidth}{0.5pt}\end{center}

\textbf{Non-square Kairo sample counts}

Current implementation chooses:

\[
s=
\left\lceil\sqrt N\right\rceil.
\]

Then uses sample index to select one cell in an \(s\times s\) grid.

For:

\[
N=16,
\]

excellent:

\[
4\times4.
\]

For:

\[
N=5,
\]

it constructs a conceptual:

\[
3\times3
\]

grid but only fills five cells.

The sample set therefore does not cover all strata symmetrically.

This can be an incorrect Monte Carlo estimator for the full pixel: equal averaging over a fixed subset of strata leaves part of the pixel with zero sampling probability. It is unbiased only under additional assumptions, or with a sampling design and weights that cover the complete pixel. Randomize the subset uniformly with correct weights, use a complete stratification, or use another full-support sampler.

Square SPP values work especially naturally with this implementation.

\begin{center}\rule{0.5\linewidth}{0.5pt}\end{center}

\textbf{Low-discrepancy sampling}

Random samples can cluster.

A low-discrepancy sequence deliberately fills space evenly.

For a sample set \(x_1,\dots,x_N\), discrepancy measures deviation between:

\[
\text{fraction of samples in a region}
\]

and:

\[
\text{region's volume}.
\]

Roughly:

\[
\boxed{
\text{low discrepancy}
=
\text{few large holes and few clusters}.
}
\]

Common families include:

\[
\text{Van der Corput},
\]

\[
\text{Halton},
\]

\[
\text{Hammersley},
\]

\[
\text{Sobol}.
\]

\begin{center}\rule{0.5\linewidth}{0.5pt}\end{center}

\textbf{Van der Corput sequence}

Take integer \(n\), write it in base 2, reverse its digits after the radix point.

For example:

\[
n=1:
1_2
\rightarrow
0.1_2
=
0.5.
\]

\[
n=2:
10_2
\rightarrow
0.01_2
=
0.25.
\]

\[
n=3:
11_2
\rightarrow
0.11_2
=
0.75.
\]

\[
n=4:
100_2
\rightarrow
0.001_2
=
0.125.
\]

So sequence begins approximately:

\[
0.5,\ 0.25,\ 0.75,\ 0.125,\ 0.625,\dots
\]

Notice how it progressively fills gaps.

That is very different from independent random clustering.

\begin{center}\rule{0.5\linewidth}{0.5pt}\end{center}

\textbf{Halton}

Halton uses radical inverse with different prime bases for each dimension.

Dimension 0:

\[
\text{base }2.
\]

Dimension 1:

\[
\text{base }3.
\]

Dimension 2:

\[
\text{base }5.
\]

and so on.

It is simple and effective at low dimensions, but higher-dimensional correlations can become problematic.

Rendering can require dozens or hundreds of dimensions because each bounce consumes multiple random values.

That is why Sobol and scrambling are common modern choices.

\begin{center}\rule{0.5\linewidth}{0.5pt}\end{center}

\textbf{Sobol sampling}

Sobol is a digital low-discrepancy sequence designed for high-dimensional integration.

One major advantage is that samples can be addressed by:

\[
(\text{sample index},\text{dimension}).
\]

This fits rendering perfectly.

Conceptually:

\begin{verbatim}
float u =
    Sobol(sampleIndex, dimension);
\end{verbatim}

Then dimension assignment remains stable even if work scheduling changes.

Scrambling, such as Owen-style scrambling, helps decorrelate sequences between pixels while preserving useful low-discrepancy structure.

\begin{center}\rule{0.5\linewidth}{0.5pt}\end{center}

\textbf{Low discrepancy is not simply ``better randomness''}

For Monte Carlo, pure independent randomness gives classic:

\[
O(N^{-1/2})
\]

root-mean-square convergence.

Quasi-Monte-Carlo low-discrepancy sampling can achieve much better behavior for sufficiently smooth, lower-effective-dimensional integrands.

But rendering integrands contain discontinuities from:

\begin{itemize}
\tightlist
\item
  visibility,
\item
  texture edges,
\item
  delta events.
\end{itemize}

So real performance depends strongly on scene/integrator structure.

You benchmark.

You do not assume Sobol beats everything in every rendering regime.

\begin{center}\rule{0.5\linewidth}{0.5pt}\end{center}

\textbf{Blue-noise sampling}

White noise has energy spread across frequencies.

Humans are especially sensitive to low-frequency noise such as:

\begin{itemize}
\tightlist
\item
  blotches,
\item
  bands,
\item
  large clusters.
\end{itemize}

Blue-noise distributions suppress low-frequency error and push energy toward higher spatial frequencies.

Visually:

\begin{verbatim}
white noise:
clusters + holes

blue noise:
even local spacing, no obvious grid
\end{verbatim}

The rendered estimator may have similar average numerical error, yet blue-noise error often looks much less objectionable.

That is crucial for low-sample real-time rendering.

\begin{center}\rule{0.5\linewidth}{0.5pt}\end{center}

\textbf{Low discrepancy versus blue noise}

Low-discrepancy sequences optimize integration uniformity over sample count/dimensions.

Blue noise optimizes spatial/frequency distribution of error.

They solve overlapping but distinct problems.

A renderer may want:

\[
\text{low discrepancy across path dimensions}
\]

while simultaneously using:

\[
\text{blue-noise scrambling across image pixels}.
\]

These ideas are complementary.

\begin{center}\rule{0.5\linewidth}{0.5pt}\end{center}

\textbf{Threaded tile rendering}

Path tracing is embarrassingly parallel at a high level:

\[
\text{each pixel/sample can usually be evaluated independently}.
\]

Instead of launching a thread per pixel, divide the image into tiles.

For example:

\[
16\times16.
\]

Each worker repeatedly obtains a tile and renders all pixels inside it.

Advantages include:

\begin{itemize}
\tightlist
\item
  good cache locality,
\item
  manageable work units,
\item
  low scheduling overhead,
\item
  no pixel-level synchronization.
\end{itemize}

Kairo already uses this architecture.

\begin{center}\rule{0.5\linewidth}{0.5pt}\end{center}

\textbf{Current Kairo tile scheduler}

The renderer computes tile count and creates:

\[
\texttt{nextTile}
\]

as:

\begin{verbatim}
std::atomic<uint32_t>
\end{verbatim}

Each worker executes:

\begin{verbatim}
tileIndex =
    nextTile.fetch_add(
        1,
        memory_order_relaxed);
\end{verbatim}

until all tiles are consumed.

This is a \textbf{dynamic work queue implemented with an atomic counter}.

It is simple and effective.

\begin{center}\rule{0.5\linewidth}{0.5pt}\end{center}

\textbf{Why dynamic work allocation matters for path tracing}

Two image tiles can have wildly different cost.

Tile A may see:

\begin{verbatim}
empty sky.
\end{verbatim}

Each path:

\[
1\text{ intersection + miss}.
\]

Tile B may see:

\begin{verbatim}
glass + mirrors + complex BVH region + many shadow rays.
\end{verbatim}

Each pixel may require many bounces and intersection queries.

Static assignment:

\begin{verbatim}
thread 0 gets fixed image quarter
thread 1 gets another quarter
\end{verbatim}

can leave one worker busy long after others finish.

Dynamic tiles allow idle workers to take more work.

That reduces tail latency.

\begin{center}\rule{0.5\linewidth}{0.5pt}\end{center}

\textbf{Tile-size tradeoff}

If tile is enormous:

\[
512\times512,
\]

there may be too few tasks.

Load imbalance increases.

If tile is tiny:

\[
1\times1,
\]

the atomic work-distribution overhead becomes comparable to rendering work.

Good tile size balances:

\[
\text{locality},
\]

\[
\text{load granularity},
\]

\[
\text{scheduling overhead}.
\]

For CPU path tracing, values around:

\[
8\times8,
\ 16\times16,
\ 32\times32
\]

are common experiment points.

The correct Kairo size should be benchmarked.

\begin{center}\rule{0.5\linewidth}{0.5pt}\end{center}

\textbf{Per-thread RNGs}

A traditional multithreaded renderer assigns each worker its own RNG state.

This eliminates locks.

Conceptually:

\begin{verbatim}
ThreadState
{
    RNG rng;
    Stats stats;
};
\end{verbatim}

But this has a reproducibility problem.

If thread scheduling changes, pixel A may be evaluated by a different worker and receive a different random stream.

Kairo's per-pixel/sample seed design avoids this.

So for Kairo, I would distinguish:

\[
\boxed{
\text{worker-local storage for speed}
}
\]

from:

\[
\boxed{
\text{pixel/sample-addressed randomness for determinism}.
}
\]

The RNG object can be local to a worker/path, but the seed should come from stable path identity.

\begin{center}\rule{0.5\linewidth}{0.5pt}\end{center}

\textbf{Thread-local statistics}

Kairo does another smart thing:

Each worker gets:

\begin{verbatim}
RenderStats& localStats =
    workerStats[workerIndex];
\end{verbatim}

and updates it without atomics.

After all workers join, statistics are reduced.

That avoids atomic increments for every ray.

Given potentially billions of ray events, atomically incrementing a global counter per intersection would be disastrous.

\begin{center}\rule{0.5\linewidth}{0.5pt}\end{center}

\textbf{But current workerStats can false-share}

\texttt{workerStats} is:

\begin{verbatim}
std::vector<RenderStats>
\end{verbatim}

with structs adjacent in memory.

Thread 0 repeatedly writes:

\begin{verbatim}
workerStats[0]
\end{verbatim}

Thread 1 repeatedly writes:

\begin{verbatim}
workerStats[1].
\end{verbatim}

Even though they are logically different objects, fields may reside on the same:

\[
64\text{-byte cache line}.
\]

If so, each core's write invalidates the other's cached line.

This is \textbf{false sharing}.

No actual data is shared logically, but cache coherence behaves as though it is.

\begin{center}\rule{0.5\linewidth}{0.5pt}\end{center}

\textbf{Cache-coherence intuition}

Imagine cache line contains:

\begin{verbatim}
[ thread0 counter | thread1 counter ]
\end{verbatim}

Core 0 modifies its half.

The coherence protocol must gain exclusive ownership of the entire line.

Core 1's copy becomes invalid.

Core 1 modifies its half.

Ownership moves back.

The line ``ping-pongs'' between cores.

Performance can collapse.

The fix is separate cache lines.

\begin{center}\rule{0.5\linewidth}{0.5pt}\end{center}

\textbf{Padded thread state}

A simple design:

\begin{verbatim}
struct alignas(64) WorkerStats
{
    RenderStats Value;

    // Padding may be added automatically because
    // of alignment / size rounding.
};
\end{verbatim}

Or use:

\begin{verbatim}
alignas(
    std::hardware_destructive_interference_size)
\end{verbatim}

when appropriate and supported.

You should benchmark before optimizing, but highly updated per-thread counters are classic false-sharing candidates.

\begin{center}\rule{0.5\linewidth}{0.5pt}\end{center}

\textbf{False sharing in the film}

Kairo writes pixels directly into one row-major:

\begin{verbatim}
std::vector<Color3f>
\end{verbatim}

but tile ownership means no two workers intentionally write the same pixel.

That avoids data races.

However, two simultaneously rendered neighboring tiles can touch different pixels located on the same cache line near their boundary.

This is possible false sharing.

Usually it is much less severe than sharing one hot counter line because each pixel is written relatively infrequently.

A tile-local accumulation buffer followed by a contiguous commit can eliminate even that if profiling says it matters.

\begin{center}\rule{0.5\linewidth}{0.5pt}\end{center}

\textbf{Kairo's current film write pattern is good}

Inside one \texttt{renderPixel()}:

\begin{verbatim}
Color3f accumulated = Black();

for samples:
    accumulated += Trace(...);

film.SetPixel(
    x,
    y,
    accumulated / spp);
\end{verbatim}

So the shared film is written:

\[
\boxed{\text{once per pixel}}
\]

rather than once per sample.

That minimizes contention and memory traffic.

Very good CPU architecture.

\begin{center}\rule{0.5\linewidth}{0.5pt}\end{center}

\textbf{Film sample accumulation}

For sample colors:

\[
X_1,\dots,X_N,
\]

the target pixel is:

\[
\bar X_N=
\frac1N
\sum_iX_i.
\]

Three common architectures are:

\begin{verbatim}
sum + sample count
online running mean
mean + M2 variance statistics
\end{verbatim}

For adaptive path tracing, the third is best because it supports convergence estimates.

\begin{center}\rule{0.5\linewidth}{0.5pt}\end{center}

\textbf{Accumulation must happen in linear HDR space}

Do \textbf{not} tone map each sample and average:

\[
\frac1N
\sum
T(X_i).
\]

Tone mapping is nonlinear, so:

\[
\boxed{
T\left(
\frac1N\sum X_i
\right)
\neq
\frac1N
\sum T(X_i).
}
\]

Correct pipeline:

\[
\boxed{
\text{sample radiance}
\rightarrow
\text{linear accumulation}
\rightarrow
\text{mean}
\rightarrow
\text{exposure}
\rightarrow
\text{tone map}
\rightarrow
\text{display encoding}.
}
\]

\begin{center}\rule{0.5\linewidth}{0.5pt}\end{center}

\textbf{Floating-point accumulation precision}

\texttt{Color3f} is three 32-bit floats.

Float has approximately:

\[
24
\]

bits of significand precision.

Relative machine epsilon around unity is roughly:

\[
10^{-7}.
\]

Now imagine accumulated sum becomes:

\[
10^6.
\]

Trying to add:

\[
10^{-3}
\]

may no longer affect the float because the representable spacing near \(10^6\) is too coarse.

That means many tiny samples can disappear numerically after very large contributions.

\begin{center}\rule{0.5\linewidth}{0.5pt}\end{center}

\textbf{Catastrophic accumulation scenario}

Suppose samples are:

\[
10^8,
1,
1,
1,
\dots
\]

A large firefly pushes sum to around:

\[
10^8.
\]

Float spacing near that magnitude is multiple units.

Subsequent contributions of order:

\[
1
\]

can lose precision.

This is another reason fireflies hurt numerics in addition to visual variance.

\begin{center}\rule{0.5\linewidth}{0.5pt}\end{center}

\textbf{Double accumulation}

A simple offline CPU renderer can maintain:

\begin{verbatim}
struct AccumColor
{
    double r;
    double g;
    double b;
};
\end{verbatim}

even if BRDF/intersection math remains float.

At final output convert to float.

This dramatically improves long-running progressive accumulation for modest CPU memory cost.

For:

\[
4K
=
8.29\text{ million pixels},
\]

RGB float accumulation:

\[
8.29M\times12
\approx99.5\text{ MB}.
\]

RGB double:

\[
8.29M\times24
\approx199\text{ MB}.
\]

That is not free, but very plausible for offline CPU rendering.

\begin{center}\rule{0.5\linewidth}{0.5pt}\end{center}

\textbf{Kahan compensated summation}

If doubles are too expensive, compensated summation tracks lost low-order terms.

For scalar:

\begin{verbatim}
void KahanAdd(
    float x,
    float& sum,
    float& compensation)
{
    const float y =
        x - compensation;

    const float t =
        sum + y;

    compensation =
        (t - sum) - y;

    sum = t;
}
\end{verbatim}

This significantly reduces error in long sums.

It costs extra arithmetic/state.

Again:

\[
\boxed{\text{benchmark and choose based on render duration/dynamic range}.}
\]

\begin{center}\rule{0.5\linewidth}{0.5pt}\end{center}

\textbf{Current Kairo internal HDR representation}

\texttt{Color3f} explicitly documents:

\[
\boxed{\text{linear RGB values may exceed }1.}
\]

That is exactly right.

The film therefore stores scene-referred HDR values before display conversion.

Current \texttt{SaveHDR()} writes this linear film directly to Radiance RGBE.

So unlike the earlier real-time Kairo shader we inspected, the ray tracer does preserve pre-tone-map HDR internally.

That is an important positive distinction.

\begin{center}\rule{0.5\linewidth}{0.5pt}\end{center}

\textbf{Radiance RGBE format}

RGBE stores:

\begin{verbatim}
R mantissa
G mantissa
B mantissa
shared exponent
\end{verbatim}

Conceptually:

\[
(R,G,B)\times2^E.
\]

The shared exponent gives huge dynamic range with compact 32 bits/pixel.

Current Kairo's \texttt{EncodeRGBE()} finds:

\[
M=\max(R,G,B),
\]

uses \texttt{frexp} to split magnitude into mantissa/exponent, and encodes RGB relative to that exponent.

This preserves values far above:

\[
1.
\]

\begin{center}\rule{0.5\linewidth}{0.5pt}\end{center}

\textbf{RGBE precision limitation}

RGBE has only 8-bit mantissa channels and one shared exponent.

Suppose:

\[
R=10000,
\]

\[
G=0.01.
\]

The exponent needed for red also determines green's precision.

Small channels relative to the maximum can lose accuracy.

For serious production compositing, OpenEXR with:

\[
16\text{-bit half}
\]

or:

\[
32\text{-bit float}
\]

channels is typically more flexible.

Kairo's RGBE support is still a major improvement over 8-bit PNG/PPM because it preserves scene dynamic range.

\begin{center}\rule{0.5\linewidth}{0.5pt}\end{center}

\textbf{HDR output versus display output}

HDR output should preserve:

\[
L_{\text{linear}}.
\]

No Reinhard.

No gamma 2.2.

No clamping to:

\[
[0,1].
\]

Display output needs:

\[
\text{exposure}
\rightarrow
\text{tone mapping}
\rightarrow
\text{transfer encoding}.
\]

Kairo respects this distinction:

\texttt{SaveHDR()} writes the linear film.

\texttt{SavePPM()} performs Reinhard and then approximate gamma encoding.

Good separation.

\begin{center}\rule{0.5\linewidth}{0.5pt}\end{center}

\textbf{Tone mapping}

A path tracer may produce:

\[
C=(12,\ 3,\ 0.8).
\]

A display cannot directly represent this in ordinary SDR.

Exposure first:

\[
\boxed{
C_e=
2^{EV}C.
}
\]

Then tone map.

Simple Reinhard:

\[
\boxed{
T(x)=\frac{x}{1+x}.
}
\]

For:

\[
x=12,
\]

\[
T(12)=12/13\approx0.923.
\]

For:

\[
x=0.8,
\]

\[
T(0.8)\approx0.444.
\]

Then encode to display space.

\begin{center}\rule{0.5\linewidth}{0.5pt}\end{center}

\textbf{Why exposure belongs before tone mapping}

Exposure is linear scaling:

\[
C_e=kC.
\]

Changing exposure represents camera/display adaptation.

If you apply exposure after tone mapping, the nonlinear curve has already compressed highlight relationships.

Correct conceptual sequence:

\[
\boxed{
C_{\text{scene}}
\rightarrow
2^{EV}
\rightarrow
T(C)
\rightarrow
\text{sRGB/display encoding}.
}
\]

\begin{center}\rule{0.5\linewidth}{0.5pt}\end{center}

\textbf{Debug mode: normals}

Current \texttt{TraceNormal()} outputs:

\[
\boxed{
C=
0.5N+0.5.
}
\]

So:

\[
N=(-1,-1,-1)
\rightarrow
(0,0,0)
\]

and:

\[
N=(1,1,1)
\rightarrow
(1,1,1).
\]

This bypasses materials and lighting.

It diagnoses:

\begin{itemize}
\tightlist
\item
  wrong winding,
\item
  flipped normals,
\item
  interpolation errors,
\item
  transform errors.
\end{itemize}

A beauty image can hide those behind shadows and textures.

Normal output cannot.

\begin{center}\rule{0.5\linewidth}{0.5pt}\end{center}

\textbf{Debug mode: albedo}

Kairo's debug path directly returns:

\begin{verbatim}
EvaluateAlbedo(scene, hit)
\end{verbatim}

This isolates:

\[
\text{material texture/factor interpretation}
\]

from lighting.

If the beauty image is wrong but albedo is correct, look downstream at:

\begin{itemize}
\tightlist
\item
  BRDF,
\item
  light,
\item
  visibility,
\item
  exposure.
\end{itemize}

If albedo itself is wrong, do not touch GGX.

Debugging should proceed in causal order.

\begin{center}\rule{0.5\linewidth}{0.5pt}\end{center}

\textbf{Debug mode: depth}

Kairo computes:

\[
d_n
=
1-
\frac{
t-d_{\text{near}}
}{
d_{\text{far}}-d_{\text{near}}
}
\]

clamped to:

\[
[0,1].
\]

Near:

\[
\text{white}.
\]

Far:

\[
\text{black}.
\]

This verifies:

\begin{itemize}
\tightlist
\item
  camera rays,
\item
  closest-hit selection,
\item
  BVH traversal,
\item
  intersection distances.
\end{itemize}

Because Kairo normalizes ray directions, hit \(t\) is actual world-space distance, which makes the depth visualization meaningful.

\begin{center}\rule{0.5\linewidth}{0.5pt}\end{center}

\begin{center}\rule{0.5\linewidth}{0.5pt}\end{center}

\hyperref[volume-reading-map]{Reading map} · \hyperref[03-visual-computing-i-ch009-point-lights-mis-and-environment-lighting]{Previous chapter} · \href{../06-GPU-Architecture-and-Real-Time-Rendering/main.pdf\#nameddest=03-visual-computing-i-ch011-bvh-traversal-and-deterministic-evidence}{Next chapter (GPU Architecture and Real-Time Rendering, Ch. 8)}
